\chapter{Standar API Sektoral: SATUSEHAT, HL7 FHIR, dan SNAP}

\section*{Tujuan Pembelajaran}
\begin{itemize}
  \item Menjelaskan alasan regulator menetapkan standar API sektoral dan dampaknya bagi organisasi.
  \item Menjelaskan konsep dasar HL7 FHIR (\textit{Health Level Seven Fast Healthcare Interoperability Resources}): \textit{resource}, profil, dan terminologi.
  \item Menilai kesenjangan integrasi organisasi terhadap standar sektor yang berlaku.
\end{itemize}

\section{Pendahuluan}

% TODO
% Cakupan: Medikara wajib mengirim data rekam medis ke SATUSEHAT dan bertukar data klaim dengan BPJS
% Kesehatan, tetapi tiap rumah sakit melakukannya dengan cara yang berbeda.

\section{Mengapa Regulator Menetapkan Standar API}

% TODO
% Cakupan: interoperabilitas sebagai infrastruktur bersama; penurunan biaya integrasi di tingkat
% nasional; persaingan yang lebih sehat karena pendatang baru dapat terhubung; standar wajib dan
% standar sukarela.

\section{HL7 FHIR: Resource, Profil, dan Terminologi}

% TODO
% Cakupan: resource sebagai unit data (Patient, Encounter, Observation); profil sebagai penyesuaian
% nasional; terminologi seperti ICD, LOINC, dan SNOMED CT; FHIR R4 (4.0.1) yang dipakai SATUSEHAT,
% R5 sebagai rilis terbit terbaru, dan R6 yang masih dalam proses ballot.
% Pendalaman teknis: membaca contoh resource Patient dan Encounter pada spesifikasi FHIR R4, lalu
% memetakan kolom data SIMRS ke elemen resource tersebut.

\section{SATUSEHAT dan Rekam Medis Elektronik}

% TODO
% Cakupan: kewajiban rekam medis elektronik dan interoperabilitas dalam Permenkes 24/2022; fungsi
% SATUSEHAT Platform; alur integrasi dari registrasi, sandbox, sampai produksi; UU 17/2023 tentang
% Kesehatan disinggung singkat.
% Kasus Medikara: lima SIMRS dengan versi berbeda yang harus mengirim data dengan standar yang sama.

\section{Integrasi dengan BPJS Kesehatan}

% TODO
% Cakupan: layanan bridging BPJS Kesehatan untuk klaim dan antrean; kredensial dan lingkungan uji;
% tantangan ketika SIMRS berbeda versi; hubungan antrean daring dengan pengalaman pasien.
% Kasus Medikara: antrean layanan yang panjang dan integrasi point-to-point tiap rumah sakit.

\section{Open Banking, Open Finance, dan SNAP}

% TODO
% Cakupan: konsep open banking dan open finance; SNAP dari Bank Indonesia beserta standar teknis
% dan keamanan, standar data, dan pedoman tata kelola; pengelolaan oleh ASPI sejak 1 September
% 2023 dan revisi standar ke versi 1.0.2 pada September 2024; layanan digital bank bersama mitra
% menurut POJK 21/2023; PSD2 dan rencana PSD3 di Uni Eropa sebagai pembanding.
% Kasus Medikara: pembayaran layanan dan klaim asuransi yang melibatkan bank atau dompet digital.

\section{Pelajaran Lintas Sektor}

% TODO
% Cakupan: perbandingan sektor kesehatan dan keuangan (wajib atau sukarela, pengelola standar,
% sertifikasi, sanksi); dampak standar sektoral terhadap strategi platform; standar sebagai
% boundary resource bersama.

\section{Aktivitas: Penilaian Kesiapan Interoperabilitas}

% TODO
% Tujuan aktivitas: mengetahui jarak antara integrasi saat ini dan standar sektor yang berlaku.
% Peran dalam tim: anggota berlatar kesehatan, hukum, atau keuangan membaca kewajiban regulasi
% beserta tenggatnya; anggota berlatar teknis membandingkan data yang dikirim saat ini dengan
% resource atau spesifikasi standar.
% Langkah: tentukan standar yang relevan (SATUSEHAT dan BPJS Kesehatan untuk Medikara, SNAP untuk
% organisasi keuangan); daftar kewajiban; nilai status tiap integrasi; tandai kesenjangan proses,
% data, dan teknologi; rangkum bersama luaran Bab 4 sampai 6.
% Luaran dan tonggak capstone: penilaian kesiapan interoperabilitas, melengkapi C2 yang ditinjau
% menjelang UTS.
% Berkas sources yang direncanakan: sources/chapter07/penilaian-kesiapan-interoperabilitas.md,
% sources/chapter07/contoh-patient-encounter-r4.json, sources/chapter07/pemetaan-simrs-fhir.csv

\section{Ringkasan}

% TODO

\section{Latihan}

\begin{enumerate}
  \item Jelaskan mengapa SATUSEHAT memakai standar terbuka seperti HL7 FHIR, bukan format buatan sendiri.
  \item Bandingkan cara Indonesia melalui SNAP (Standar Nasional Open API Pembayaran) dan Uni Eropa melalui PSD2 (\textit{Payment Services Directive 2}) mendorong \textit{open banking}. Sebutkan satu persamaan dan satu perbedaan.
  \item Misalkan data kunjungan pasien dari dua rumah sakit Medikara belum terkirim ke SATUSEHAT. Uraikan tiga kemungkinan penyebab dari sisi proses, data, dan teknologi.
\end{enumerate}
